 \documentclass[xcolor=dvipsnames,t,8pt,fleqn]{beamer}
%\documentclass[aspectratio=1610,8pt]{beamer}

%Empfehlungen aus der ct 9/2007
%\usepackage{beamerthemesplit}
%\usepackage[T1]{fontenc}
%\usepackage{beamer}

\usepackage{mathptmx}
\usepackage{helvet}
\usepackage{eurosym}

% Layout-Vorschl�ge - ausw�hlen!
% Siehe hierzu das Handbuch beameruserguide.pdf
% 1. Presentation Theme
\usetheme{Goettingen}
% 2. Color Theme
\usecolortheme{sidebartab}
%\usecolortheme{dove}
%\usecolortheme{lily}  % Inner Color Theme
% 3. Font Theme
\usefonttheme{professionalfonts}
% 4. Inner Theme
\useinnertheme{circles}
% 5. Outer Theme
%\useoutertheme{sidebar}
%
%\usepackage{extsizes} % Da 8pt in der Pr�sentation

\usepackage{graphicx}
%amsthm wird f�r styles bei Satz- und Definitionsnumerierung ben�tigt
%bbm wird f�r die Symbole f�r reelle bzw. nat. Zahlen ben�tigt
\usepackage{fancyhdr, amssymb, amsmath, amsthm, bbm} 
\usepackage[ngerman]{babel}
%\usepackage{ifpdf}

%Umlaute
\usepackage[latin1]{inputenc} 
\usepackage[T1]{fontenc}
\usepackage{lmodern}

% evtl. f�r Code
\usepackage{listings}
% links to URLs
\usepackage{hyperref}
%
\usepackage{tikz}
\usetikzlibrary{matrix}
\usetikzlibrary{positioning}
\usetikzlibrary{backgrounds}
%\usetikzlibrary{intersections}
%\usetikzlibrary{calc}
%
\usepackage{pgfplots}
%
\usepackage{venndiagram}
%
\pgfplotsset{compat=1.15}
%
%\usepackage{showframe}
%
\setbeamercovered{transparent}
%bewirkt, dass schrittweise Einblendungen zuvor schon transparent angezeigt werden
\setbeamertemplate{footline}[page number]

\title[Data Science]{Data Science}
%[Kurztitel] - erscheint auf allen Folien
%{Titel} - erscheint auf der Titelseite (wird mit \titlepage dargestellt, s.u.)
\subtitle{Chapter 7: Model Evaluation}
%Untertitel - kommt auch auf die Titelseite
\author{S.~Steidel \& F.~Wick}
%Autor kommt auf den Titel und - je nach theme - auf die �bersichtsleiste
\date{\today}
%Datum kommt auf den Titel - hier das aktuelle Datum bei Erstellung der Folien

%%%%%%%\DeclareMathOperator{\Arg}{Arg}
%%%%%%%\DeclareMathOperator{\e}{e}
%%%%%%%\DeclareMathOperator{\Inv}{Inv}
%%%%%%%\DeclareMathOperator{\id}{id}
%%%%%%%\DeclareMathOperator{\rot}{rot}
%%%%%%%\DeclareMathOperator{\grad}{grad}
%%%%%%%\DeclareMathOperator{\Div}{div}
%%%%%%%\DeclareMathOperator{\Kern}{Kern}

\usepackage{subfig}

%\usepackage{amsmath}
\usepackage[intlimits]{empheq}
\usepackage{amsthm}
\usepackage{amsfonts}
%\usepackage{latexsymb}
\usepackage{pifont}
%\usepackage[amsmath]{ntheorem}
\usepackage{trsym}
\usepackage{trfsigns}
\usepackage{wasysym}
%
\usepackage[derivedinbase, derived, thinspace, thinqspace, squaren, binary]{SIunits} 
\usepackage[intlimits]{empheq}
%
\input{Stochastik-Makros.tex}
%
%%%%%%%%%%%%%%%%%%%%%%%%%%%%%%%%%%%%%%%%%%%%%%%%%%%%%%%%%%%%%%%%%%%%%%%%%%%%%%%%%%%%%%%%%%%%%%%%%%%%%%
%%%%%%%%%%%%%%%%%%%%%%%%%%%%%%%%%%%%%%%%%%%%%%%%%%%%%%%%%%%%%%%%%%%%%%%%%%%%%%%%%%%%%%%%%%%%%%%%%%%%%%
%%%%%%%%%%%%%%%%%%%%%%%%%%%%%%%%%%%%%%%%%%%%%%%%%%%%%%%%%%%%%%%%%%%%%%%%%%%%%%%%%%%%%%%%%%%%%%%%%%%%%%
%
\begin{document}
%
%%%%%%%\newcommand{\R}{\mathbbm{R}}
%%%%%%%\newcommand{\N}{\mathbbm{N}}
%%%%%%%\newcommand{\Z}{\mathbbm{Z}}
%%%%%%%\newcommand{\Q}{\mathbbm{Q}}
%%%%%%%\newcommand{\C}{\mathbbm{C}}
%%%%%%%\newcommand{\disp}{\displaystyle}
%%%%%%%\newcommand{\bs}{\boldsymbol}
%%%%%%%\let\oldhat\hat
%%%%%%%\renewcommand{\vec}[1]{\boldsymbol{#1}}
%%%%%%%\renewcommand{\hat}[1]{\oldhat{\boldsymbol{#1}}}
%
%%%%%%%%%%%%%%%%%%%%%%%%%%%%%%%%%%%%%%%%%%%%%%%%%%%%%%%%%%%%%%%%%%%%%%%%%%%%%%%%%%%%%%%%%%%%%%%%%%%%%%
%%%%%%%%%%%%%%%%%%%%%%%%%%%%%%%%%%%%%%%%%%%%%%%%%%%%%%%%%%%%%%%%%%%%%%%%%%%%%%%%%%%%%%%%%%%%%%%%%%%%%%
\begin{frame}[label=Titelrahmen]
	\titlepage
\end{frame}
%
%%%\begin{frame}
	%%%\frametitle{Standardsignale}
%%%\end{frame}
%%%%%%%%%%%%%%%%%%%%%%%%%%%%%%%%%%%%%%%%%%%%%%%%%%%%%%%%%%%%%%%%%%%%%%%%%%%%%%%%%%%%%%%%%%%%%%%%%%%%%%
%%%%%%%%%%%%%%%%%%%%%%%%%%%%%%%%%%%%%%%%%%%%%%%%%%%%%%%%%%%%%%%%%%%%%%%%%%%%%%%%%%%%%%%%%%%%%%%%%%%%%%
%%%%%%%%%%%%%%%%%%%%%%%%%%%%%%%%%%%%%%%%%%%%%%%%%%%%%%%%%%%%%%%%%%%%%%%%%%%%%%%%%%%%%%%%%%%%%%%%%%%%%%
\setcounter{section}{7}
%%%%%%%%%%%%%%%%%%%%%%%%%%%%%%%%%%%%%%%%%%%%%%%%%%%%%%%%%%%%%%%%%%%%%%%%%%%%%%%%%%%%%%%%%%%%%%%%%%%%%%
%%%%%%%%%%%%%%%%%%%%%%%%%%%%%%%%%%%%%%%%%%%%%%%%%%%%%%%%%%%%%%%%%%%%%%%%%%%%%%%%%%%%%%%%%%%%%%%%%%%%%%


\section{Hypothesis Testing / Parameter Testing}

\subsection{Introduction}

\begin{frame}%[allowframebreaks,allowdisplaybreaks]
	\frametitle<presentation>{Introduction}
	%
	\begin{itemize}
		\item[\ding{43}] So far: \alert{parameter} and \alert{interval estimation} for a parameter $\theta$ of a \alert{probabilistic model} or \alert{distribution} from an empirical sample
		\item[\ding{45}] Now: Making a justified decision about the \emph{validity} or \emph{invalidity} of a \alert{hypothesis} about the parameters of a \alert{probabilistic model} or distribution based on available observations ($=$ sample(s)): \alert{parameter testing}
	\end{itemize}
	%
	\par\smallskip
	%
	\begin{itemize}
		\item Parameter tests are a subset of \alert{statistical tests}. Statistical testing means using the information gained (computed) from a sample to make a decision about a \alert{hypothesis}: \alert{hypothesis test} 
		
		\item A \alert{hypothesis} is an (unproven) assumption about:
		\begin{itemize}
			\item the distribution, or
			\item specific parameters of a population.
		\end{itemize}
		
		\item \alert{Hypotheses} may arise from:
		\begin{itemize}
			\item previous statistical observations,
			\item theoretical considerations, or
			\item informed conjectures.
		\end{itemize}
		
		\item Hypothesis tests do \textit{not} decide or \textit{cannot} prove whether the hypothesis is \textit{false} or \textit{true}. As a result of the hypothesis test, the hypothesis is either \alert{retained} ("`accepted"') or \alert{rejected} ("`discarded"'). 
		
		\item Important: Control the probabilities of errors:
		\begin{itemize}
			\item Rejecting a true hypothesis,
			\item Retaining a false hypothesis.
		\end{itemize} 
	\end{itemize}
	%
\end{frame}

%\begin{frame}{Statistical Inference: Estimation and Hypothesis Testing}
%	
%	\begin{itemize}
%		\item Estimation methods and hypothesis testing are key applications of sampling theory.
%		
%		\item In hypothesis testing, information obtained from a sample is used to make a decision about a hypothesis.
%		
%		\item A \textbf{hypothesis} is an assumption about:
%		\begin{itemize}
%			\item the distribution, or
%			\item specific parameters of a population.
%		\end{itemize}
%		
%		\item Hypotheses may arise from:
%		\begin{itemize}
%			\item previous statistical observations,
%			\item theoretical considerations,
%			\item the principle of insufficient reason, or
%			\item informed conjectures.
%		\end{itemize}
%		
%		\item Similar to estimation, hypothesis testing cannot provide certainty:
%		\begin{itemize}
%			\item A hypothesis cannot be proven true or false with absolute certainty.
%			\item The outcome is to \textbf{retain} or \textbf{reject} the hypothesis.
%		\end{itemize}
%		
%		\item Important: Control the probabilities of errors:
%		\begin{itemize}
%			\item Rejecting a true hypothesis,
%			\item Retaining a false hypothesis.
%		\end{itemize}
%		
%	\end{itemize}
%	
%\end{frame}


\subsection{Basic principle}
\subsubsection{Null and alternative hypothesis}
%%%%%%%%%%%%%%%%%%%%%%%%%%%%%%%%%%%%%%%%%%%%%%%%%%%%%%%%%%%%%%%%%%%%%%%%%%%%%%%%%%%%%%%%%%%%%%%%%%%%%%
\begin{frame}[allowframebreaks,allowdisplaybreaks]
	\frametitle<presentation>{Null and alternative hypotheses}
	%
	The \alert{null hypothesis} or \alert{initial hypothesis} ${\mathrm{ H}}_0$ is an assumption about parameters $\theta$ of a distribution, shown here for a numerical value of a parameter, for example
	%
	\begin{equation*}
		\label{eq:15-1}
		{\mathrm{ H}}_0: \theta \leq \theta_0 \quad\text{oder}\quad {\mathrm{ H}}_0: \theta = \theta_0 \quad\text{oder}\quad {\mathrm{ H}}_0: \theta \geq \theta_0
	\end{equation*}
	%
	The null hypothesis is regarded as correct until \emph{sufficient evidence} to the contrary has been provided.
	\par\bigskip
	%
	The \alert{alternative hypothesis} ${\mathrm{ H}}_1$ is a complementary assumption to ${\mathrm{ H}}_0$: 
	\begin{equation*}
		{\mathrm{ H}}_1 = {\overline{\mathrm{ H}}}_0
	\end{equation*}
	%
	\alert{Two-sided questions} test whether the parameter of interest $\theta$ does not have the value $\theta_0$. The hypotheses are
	%
	\begin{equation*}
		\label{eq:15-1a}
		{\mathrm{ H}}_0: \theta = \theta_0 \quad\text{against}\quad {\mathrm{ H}}_1: \theta \neq \theta_0 \quad\text{"`two-sided test"'}
	\end{equation*}
	%
	\alert{One-sided questions} test whether the parameter of interest $\theta$ is smaller or larger than the value $\theta_0$, and the possible hypotheses are
	\begin{eqnarray*}
		\label{eq:15-1b}
		{\mathrm{ H}}_0: \theta \leq \theta_0 &\quad\text{against}\quad& {\mathrm{ H}}_1: \theta > \theta_0 \quad\text{"`upper-tailed test"'}\\
		\label{eq:15-1c}
		{\mathrm{ H}}_0: \theta \geq \theta_0 &\quad\text{against}\quad& {\mathrm{ H}}_1: \theta < \theta_0 \quad\text{"`lower-tailed test"'}
	\end{eqnarray*} 
	%
	\newpage
	%
	When formulating hypotheses, the following must be strictly observed:
	\begin{itemize}
		\item The \alert{null hypothesis} and all other logically possible statements must be mutually exclusive, i.e.\,the \alert{null hypothesis} must be formulated so that there is only one other statement: the \alert{alternative hypothesis}.
		\item The \alert{null hypothesis} should \textbf{not} be the statement that is to be \textbf{confirmed} by the test procedure and the sample result, but rather the \textbf{logical opposite statement}. 
		\par\medskip
		\emph{Example: A political scientist might begin an empirical study with the null hypothesis that voter support for a chancellor candidate is the same among men and women.}
		\par\medskip
		This is due to the \alert{falsification principle}:
		\begin{itemize}
			\item \textbf{Refuting a statement generally yields more valuable information than confirming it}.
			\item If a statement is confirmed by the sample result, this does not mean that it is correct, but only that the observed data are not sufficient to refute it.
		\end{itemize}  
		\par\smallskip
		Therefore, instead of speaking of accepting the \alert{null hypothesis}, it is better to speak of "`retaining"' or even "`not rejecting"' it.
	\end{itemize}
\end{frame}
%

\subsubsection{Type 1 and type 2 errors}
%%%%%%%%%%%%%%%%%%%%%%%%%%%%%%%%%%%%%%%%%%%%%%%%%%%%%%%%%%%%%%%%%%%%%%%%%%%%%%%%%%%%%%%%%%%%%%%%%%%%%%
\begin{frame}%[allowframebreaks,allowdisplaybreaks]
	\frametitle<presentation>{Type 1 and type 2 errors}
	%
	When testing hypotheses, two errors can occur:
	\par\bigskip
	\begin{center}
		\begin{tabular}{r|c|c} %\firsthline
			& ${\mathrm{ H}}_0$ \textbf{accept} & ${\mathrm{ H}}_0$ \textbf{reject} \\ \hline
			${\mathrm{ H}}_0$ is \textbf{true} & correct decision               & \alert{type 1 error}         \\ \hline
			${\mathrm{ H}}_0$ is \textbf{false}  & \alert{type 2 error}               & correct decision \\ \hline
		\end{tabular}
	\end{center}
	\par\bigskip
	\begin{enumerate}
		\item \alert{type 1 error}: The null hypothesis is rejected even though ${\mathrm{ H}}_0$ is true. The probability of type 1 error is
		\begin{equation*}
			\nonumber
			P\left( {\mathrm{ H}}_0\text{ reject} \mid {\mathrm{ H}}_0\text{ is true} \right) = \alpha 
		\end{equation*}
		\item \alert{type 2 error}: The null hypothesis is not rejected even though ${\mathrm{ H}}_0$ is false. The probability of type 2 error is
		\begin{equation*}
			\nonumber
			P\left( {\mathrm{ H}}_0\text{ accept} \mid {\mathrm{ H}}_0\text{ is false} \right) = \beta 
		\end{equation*}
	\end{enumerate}
	%
	\par\bigskip
	Usually, the type 1 error is specified, because one fears the risk of type 1 error more; typical values are $\alpha = 5\%$ and $\alpha = 1\%$. 
	\par\smallskip
	The type 2 error must then be computed from this specification (usually very complicated), because $\beta = 1 - \alpha$ does \textit{not} hold!
	%
\end{frame}

%
\subsubsection{Decision}

\begin{frame}[allowframebreaks,allowdisplaybreaks]
	\frametitle<presentation>{Decision}
	%
	After the hypotheses ${\mathrm{ H}}_0$ and ${\mathrm{ H}}_1$ have been formulated, the distribution $F_{\Xi}(\xi)$ of a suitable \alert{test statistic $\Xi$} is derived under the \alert{condition that the null hypothesis is true}.
	\par\medskip
	Then \alert{disjoint acceptance and rejection regions} (usually intervals) are defined so that the \alert{type 1 error is at most $\alpha$}:
	\begin{equation*}
		\nonumber
		P\left( \xi\in\text{reject region} \mid {\mathrm{ H}}_0\text{ is true} \right) = \alpha 
	\end{equation*} 
	The probability of a type 1 error, $\alpha$, is also called the \alert{significance level} of the test.
	\par\medskip
	Next, the \alert{concrete test statistic $\xi$} (as a realization of $\Xi$) is computed from the sample and checked to see whether it falls into the acceptance or rejection region. \alert{If it falls into the acceptance region, the null hypothesis ${\mathrm{ H}}_0$ is retained; otherwise, it is rejected in favor of ${\mathrm{ H}}_1$.} 
	\par\bigskip
	Even though the theoretical background of hypothesis tests is anything but trivial, their practical application is comparatively simple \smiley.
	
	\framebreak
	
	Acceptance and Rejection region for two-sided and one-sided hypothesis tests [Schira, Statistische Methoden der VWL und BWL, 2021]:
	\par%\smallskip
	\begin{center}
		\includegraphics[height=0.65\textwidth]{Bilder/AcceptRejectRegion.png}
	\end{center}
	\footnotesize{\textbf{Legend:} two-sided test = zweiseitiger Test; upper-sided test = oberseitiger Test; lower-sided test = unterseitiger Test, acceptance region = Annahmebereich; rejection region = Verwerfungsbereich}
	
\end{frame}


\subsection{Specific variants}

\subsubsection{Two-sided test}
\begin{frame}%[allowframebreaks,allowdisplaybreaks]
	\frametitle<presentation>{Two-sided test}
	%
	\begin{align*}
		\text{\alert{Hypotheses:} }          & {\mathrm{ H}}_0: \theta = \theta_0 \quad\text{against}\quad {\mathrm{ H}}_1: \theta \neq \theta_0 \\
		\text{\alert{Acceptance region:} }      & W_0 = [\xi_{\mathrm{ c,u}}, \xi_{\mathrm{ c,o}}] \\
		\text{\alert{Rejection regions:} } & W_1 = (-\infty,\xi_{\mathrm{ c,u}}) \cup (\xi_{\mathrm{ c,o}}, \infty)
	\end{align*}
	The corresponding probabilities are
	{\usebeamerfont{Gr07_00}
		\begin{align*}
			P\left(\Xi \in W_1 \right) &= \underbrace{P\left( \Xi < \xi_{\mathrm{ c,u}} \right)}_{\stackrel{!}{=} \alpha/2} + \underbrace{P\left( \Xi > \xi_{\mathrm{ c,o}} \right)}_{\stackrel{!}{=} \alpha/2} = \underbrace{F_{\Xi}(\xi_{\mathrm{ c,u}})}_{= \alpha/2} +  \underbrace{1 -F_{\Xi}(\xi_{\mathrm{ c,o}})}_{= \alpha/2} = \alpha \\
			P\left(\Xi \in W_0 \right) &= P\left( \xi_{\mathrm{ c,u}} \leq \Xi \leq \xi_{\mathrm{ c,o}}  \right)  = F_{\Xi}(\xi_{\mathrm{ c,o}})  - F_{\Xi}(\xi_{\mathrm{ c,u}}) =1-\alpha
		\end{align*}
	}where the \alert{significance level $\alpha$} is split symmetrically. The values $\xi_{\mathrm{ c,u}}$ and $\xi_{\mathrm{ c,o}}$ are called the "`critical lower"' and "`critical upper"' values, respectively, and they follow from the \alert{$\alpha/2$ and $(1-\alpha/2)$ quantiles} of the distribution $F_{\Xi}(\xi)$ as
	\begin{equation*}
		\xi_{\mathrm{ c,u}} = \xi_{[\alpha/2]} \text{ and } \xi_{\mathrm{ c,o}} = \xi_{[1-\alpha/2]}
	\end{equation*}
	Decision rule:
	\begin{equation*}
		\xi < \xi_{[\alpha/2]} \text{ or } \xi > \xi_{[1-\alpha/2]} \quad \Longrightarrow\quad\text{reject ${\mathrm{ H}}_0$!}
	\end{equation*}
	Decision rule (\textit{symmetric} distribution $F_{\Xi}(\xi)$):
	\begin{equation*}
		\left| \xi \right|  > \xi_{[1-\alpha/2]} \quad \Longrightarrow\quad\text{reject ${\mathrm{ H}}_0$!}
	\end{equation*}
	
\end{frame}

\begin{frame}{Example: two-sided test}
	
	\textbf{Problem:} \\
	In a student pub, beer glasses are supposed to contain $0.4\,l$. \\
	A sample of size $n=50$ yields:
	\begin{itemize}
		\item Sample mean: $\bar{x} = 0.38\,l$
		\item Sample variance: $s^2 = 0.0064\,l^2$
	\end{itemize}
%	$$\text{Sample mean:} \; \bar{x} = 0.38\,l;
%	\quad \text{Sample variance:} \; s^2 = 0.0064\,l^2$$
	
	\vspace{0.5em}
	
	\textbf{Step 1: Hypotheses}
	\[
	{\mathrm{ H}}_0: \mu = 0.4\,l \quad \text{vs.} \quad {\mathrm{ H}}_1: \mu \neq 0.4\,l
	\]
	
	\textbf{Step 2: Standard Deviation}
	\[
	\hat{\sigma}_{\bar{X}} = \sqrt{\tfrac{s^2}{n}} = \sqrt{\tfrac{0.0064\,l^2}{50}} \approx 0.01143\,l
	\]
	
	\textbf{Step 3: Test Statistic}
	\[
	z = \frac{\bar{x} - \mu_0}{\hat{\sigma}_{\bar{X}}}
	= \frac{0.38 - 0.4}{0.01143} \approx -1.75
	\]
	
	\textbf{Step 4: Significance Level} ($\alpha = 0.05$) \textbf{\& Critical Value}
	\[
	z_{[1-0.025]} = z_{[0.975]} = 1.96
	\]
	
	\textbf{Step 5: Decision}
	\[
	|z| = 1.75 < 1.96 \quad \Longrightarrow \quad \text{retain or do not reject } {\mathrm{ H}}_0
	\]
	
	\vspace{0.5em}
	\textbf{Conclusion:} \\
	There is no significant evidence that the mean filling amount differs from $0.4\,l$.
	
\end{frame}


\subsubsection{Upper-sided test}
\begin{frame}%[allowframebreaks,allowdisplaybreaks]
	\frametitle<presentation>{Upper-sided test}
	\begin{align*}
		\text{\alert{Hypotheses:} }         & {\mathrm{ H}}_0: \theta \leq \theta_0 \quad\text{against}\quad {\mathrm{ H}}_1: \theta > \theta_0 \\
		\text{\alert{Acceptance region:} }     & W_0 = (-\infty, \xi_{\mathrm{ c}}] \\
		\text{\alert{Rejection region:} } & W_1 = (\xi_{\mathrm{ c}}, \infty) 
	\end{align*}
	The corresponding probabilities are
	\begin{align*}
		P\left(\Xi \in W_1 \right) &= \underbrace{P\left( \Xi > \xi_{\mathrm{ c}} \right)}_{\stackrel{!}{=} \alpha} = \underbrace{1 -F_{\Xi}(\xi_{\mathrm{ c}})}_{= \alpha} = \alpha \\
		P\left(\Xi \in W_0 \right) &= P\left( \Xi \leq \xi_{\mathrm{ c}}  \right) =  F_{\Xi}(\xi_{\mathrm{ c}})  =1-\alpha
	\end{align*}
	where the \alert{significance level $\alpha$} is placed on the \textit{right} (= upper) side. The value $\xi_{\mathrm{ c}}$ is called the "`critical"' value, and it follows from the \alert{$(1-\alpha)$ quantile} of the distribution $F_{\Xi}(\xi)$ as
	\begin{equation*}
		\xi_{\mathrm{ c}} = \xi_{[1-\alpha]}
	\end{equation*}
	Decision rule:
	\begin{equation*}
		\xi > \xi_{[1-\alpha]} \quad \Longrightarrow\quad\text{reject ${\mathrm{ H}}_0$!}
	\end{equation*}
\end{frame}

\begin{frame}{Example: upper-sided test}
	
	\textbf{Problem:} \\
	A large trucking company pays the traffic fines its $5360$ drivers on German highways and has allocated a budget of $1000$\,\euro\,per driver/year. A random sample of $n=50$ drivers reveals that the average cost per driver last year was $\bar{x} = 1020$\,\euro. Based on many years of experience, the standard deviation has remained fairly constant at $\sigma = 60$\,\euro.
	
	\vspace{0.5em}
	
	\textbf{Step 1: Hypotheses}
	\[
	{\mathrm{ H}}_0: \mu \leq 1000\,\text{\euro} \quad \text{vs.} \quad {\mathrm{ H}}_1: \mu > 1000\,\text{\euro}
	\]
	
	\textbf{Step 2: Standard Deviation}
	\[
	\sigma_{\bar{X}} = \frac{\sigma}{\sqrt{n}} = \frac{60\,\text{\euro}}{\sqrt{50}} \approx 8.485\,\text{\euro}
	\]
	
	\textbf{Step 3: Test Statistic}
	\[
	z = \frac{\bar{x} - \mu_0}{\sigma_{\bar{X}}}
	= \frac{1020\,\text{\euro} - 1000\,\text{\euro}}{8.49\,\text{\euro}} \approx 2.3571
	\]
	
	\textbf{Step 4: Significance Level} ($\alpha = 0.05$) \textbf{\& Critical Value}
	\[
	z_{[1-0.05]} = z_{[0.95]} = 1.645
	\]
	
	\textbf{Step 5: Decision}
	\[
	z = 2.3571 > 1.645 \quad\Longrightarrow\quad \text{reject } {\mathrm{ H}}_0
	\]
	
	\vspace{0.5em}
	
	\textbf{Conclusion:} \\
	The average cost per driver is significantly above $1000\,\text{\euro}$.
	
\end{frame}


\subsubsection{Lower-sided test}
\begin{frame}%[allowframebreaks,allowdisplaybreaks]
	\frametitle<presentation>{Lower-sided test}
	\begin{align*}
		\text{\alert{Hypotheses:} }         & {\mathrm{ H}}_0: \theta \geq \theta_0 \quad\text{against}\quad {\mathrm{ H}}_1: \theta < \theta_0 \\
		\text{\alert{Acceptance region:} }     & W_0 = [\xi_{\mathrm{ c}}, \infty) \\
		\text{\alert{Rejection region:} } & W_1 = (-\infty, \xi_{\mathrm{ c}})
	\end{align*}
	The corresponding probabilities are
	\begin{align*}
		P\left(\Xi \in W_1 \right) &= \underbrace{P\left( \Xi \leq \xi_{\mathrm{ c}} \right)}_{\stackrel{!}{=} \alpha} = \underbrace{F_{\Xi}(\xi_{\mathrm{ c}})}_{= \alpha} = \alpha \\
		\nonumber
		P\left(\Xi \in W_0 \right) &= P\left( \Xi > \xi_{\mathrm{ c}}  \right)  =  1-F_{\Xi}(\xi_{\mathrm{ c}})  =1-\alpha
	\end{align*}
	where the \alert{significance level $\alpha$} is placed on the \textit{left} (= lower) side. The value $\xi_{\mathrm{ c}}$ is called the "`critical"' value, and it follows from the \alert{$\alpha$ quantile} of the distribution $F_{\Xi}(\xi)$ as
	\begin{equation*}
		\xi_{\mathrm{ c}} = \xi_{[\alpha]}
	\end{equation*}
	Decision rule:
	\begin{equation*}
		\xi < \xi_{[\alpha]} \quad \Longrightarrow\quad\text{reject ${\mathrm{ H}}_0$!}
	\end{equation*}
\end{frame}

\begin{frame}{Example: lower-sided test}
	
	\textbf{Problem:} \\
	In the previous example w.r.t.\,beer glasses, a \emph{one-sided test} would be more appropriate, since no one is going to complain if there is too much beer in the glass. 
	The	null hypothesis would therefore need to be modified.
	
	\vspace{0.5em}
	
	\textbf{Step 1: Hypotheses}
	\[
	{\mathrm{ H}}_0: \mu \geq 0.4\,l \quad \text{vs.} \quad {\mathrm{ H}}_1: \mu < 0.4\,l
	\]
	
	\textbf{Step 2: Standard Deviation}
	\[
	\hat{\sigma}_{\bar{X}} = \sqrt{\tfrac{s^2}{n}} = \sqrt{\tfrac{0.0064\,l^2}{50}} \approx 0.01143\,l
	\]
	
	\textbf{Step 3: Test Statistic}
	\[
	z = \frac{0.38 - 0.4}{0.01143} \approx -1.75
	\]
	
	\textbf{Step 4: Significance Level} ($\alpha = 0.05$) \textbf{\& Critical Value}
	\[
	z_{[0.05]} = -1.645
	\]
	
	\textbf{Step 5: Decision}
	\[
	z < -1.645 \quad\Longrightarrow\quad\text{reject } {\mathrm{ H}}_0
	\]
	
	\vspace{0.5em}
	
	\textbf{Conclusion:} \\
	Evidence that the mean filling is \textbf{below} $0.4\,l$.
	
\end{frame}


\subsubsection{Test for equality}
%
\begin{frame}%[allowframebreaks,allowdisplaybreaks]
	\frametitle<presentation>{Test for equality}
	The known principle for two-sided and one-sided testing a parameter $\theta$ can directly be transferred:
	\par\bigskip
	Two different, independently drawn samples with sizes $n_1$ and $n_2$ are considered. The \textbf{test for equality} checks whether the two populations have the same numerical value of a certain distribution parameter.
	\par\bigskip
	\alert{Two-sided questions} test whether the two parameters in question, $\theta_1$ and $\theta_2$, are equal:
	%
	\begin{equation*}
		\label{eq:15-1a}
		{\mathrm{ H}}_0: \theta_1 = \theta_2 \qquad\text{against}\qquad {\mathrm{ H}}_1: \theta_1 \neq \theta_2 \quad\text{"`two-sided test"'}
	\end{equation*}
	%
	\alert{One-sided questions} test whether the parameter in question $\theta_1$ is smaller or greater than $\theta_2$. The possible hypotheses are
	\begin{eqnarray*}
		\label{eq:15-1b}
		{\mathrm{ H}}_0: \theta_1 \leq \theta_2 &\quad\text{against}\quad& {\mathrm{ H}}_1: \theta_1 > \theta_2 \quad\text{"`upper-tailed test"'}\\
		\label{eq:15-1c}
		{\mathrm{ H}}_0: \theta_1 \geq \theta_2 &\quad\text{against}\quad& {\mathrm{ H}}_1: \theta_1 < \theta_2 \quad\text{"`lower-tailed test"'}
	\end{eqnarray*} 
	%\par\smallskip
	Strictly speaking, ${\mathrm{ H}}_0: \theta_1 = \theta_2$, since $\theta_1 = \theta_2$ is assumed for the test distribution; therefore, $\alpha$ may also be smaller (conservative estimate).
	\par\bigskip
	The procedure for making the decision is analogous to the known principle for testing a parameter $\theta$.
	%
\end{frame}

\begin{frame}{Example: test for equality}
	
	\textbf{Problem:} \\
	Lamps from two companies $X$ and $Y$ are being tested for their lifetime.\\
	Samples are drawn with $n_X = 64$, $\bar{x} = 1512\,\text{h}$, $s_X = 12\,\text{h}$ and $n_Y = 58$, $\bar{y} = 1504\,\text{h}$, $s_Y = 15\,\text{h}$.
%	\[
%	n_X = 64,\; \bar{x} = 1512\,\text{h},\; s_X = 12\,\text{h} 
%	\quad \text{and} \quad
%	n_Y = 58,\; \bar{y} = 1504\,\text{h},\; s_Y = 15\,\text{h}
%	\]	
	May one conclude that company $X$ supplies lamps with a longer lifetime?
	
	\vspace{0.5em}
	
	\textbf{Step 1: Hypotheses}
	\[
	{\mathrm{ H}}_0: \mu_X \leq \mu_Y \quad \text{vs.} \quad {\mathrm{ H}}_1: \mu_X > \mu_Y
	\]
	
	\textbf{Step 2: Standard Deviation}
	\[
	\hat{\sigma}_{\Delta} = \sqrt{
		\frac{n_X\cdot s_X^2 + n_Y\cdot s_Y^2}{n_X+n_Y-2}}
		\cdot \sqrt{\frac{n_X+n_Y}{n_X \cdot n_Y}
	}
	\approx 2.47
	\]
	
	\textbf{Step 3: Test Statistic}
	\[
	t = \frac{\bar{x} - \bar{y}}{\hat{\sigma}_{\Delta}}
	= \frac{1512 - 1504}{2.47} \approx 3.24
	\]
	
	\textbf{Step 4: Significance Level} ($\alpha=0.025$) \textbf{\& Critical Value}
	\[
	t_{n_X+n_Y-2,[1-\alpha]} = t_{120,\,0.975} \approx 1.98 \quad\text{(\emph{Student-t-distribution})}
	\]
	
	\textbf{Step 5: Decision}
	\[
	t = 3.24 > 1.98 \quad\Longrightarrow\quad \text{reject } {\mathrm{ H}}_0
	\]
	
	\vspace{0.5em}
	
	\textbf{Conclusion:} \\
	Company $X$ produces lamps with significantly longer lifetimes.
	
\end{frame}


\subsection{Summary}
\begin{frame}%[allowframebreaks,allowdisplaybreaks]
	\frametitle<presentation>{Summary \& illustration}
	%
	\begin{figure}% trim=left bottom right top, clip
		\includegraphics[trim=1.2cm 0cm 0cm 0cm, keepaspectratio, width=1.0\linewidth]{Bilder/Hypothesen_cm_embed.pdf}%
	\end{figure}
	%
	\par\smallskip
	The formal procedure for carrying out a hypothesis test always consists of the following steps:
	\begin{enumerate}
		\item Formulate the hypotheses ${\mathrm{ H}}_0$ and ${\mathrm{ H}}_1$
		\item From the sample: compute
		\begin{enumerate}
			\item[a)] required values (based on point estimates)
			\item[b)] test statistic $\xi$
		\end{enumerate}
		\item Specify $\alpha$ and determine the critical value(s) (quantile(s) of the test distribution $F_{\Xi}(\xi)$)
		\item Make the test decision by comparing the critical value(s) and the test statistic $\xi$
	\end{enumerate}
\end{frame}


\subsection{$P$-value \& significance}
\begin{frame}[allowframebreaks,allowdisplaybreaks]
	\frametitle<presentation>{$P$-value \& significance}
	%
	\begin{itemize}
		\item[\ding{43}] So far: Specification of $\alpha$ and no consideration of the distance between the computed test statistic $\xi$ and the critical value relevant for the test decision, which is defined by $\alpha$.
		\item[\ding{45}] Now: Specification of the probability \alert{$\tilde{\alpha}$} with which the computed test statistic $\xi$, under a true null hypothesis, would lie exactly on the boundary(ies) between acceptance and rejection region(s): \alert{$P$-value}. The smaller this value, the more \alert{significant} the rejection of the null hypothesis is, i.e. the more willing one is to accept the alternative hypothesis.
	\end{itemize}
	%
	\par\bigskip
	\begin{alertblock}{Definition of $\tilde{\alpha}$ ($P$-value)}
		\begin{itemize}
			\item Upper-sided test:
			\begin{equation*}
				\tilde{\alpha} = P(\Xi \geq \xi) = 1-F_{\Xi}(\xi)
			\end{equation*}
			\item Lower-sided test:
			\begin{equation*}
				\tilde{\alpha} = P(\Xi \leq \xi)  = F_{\Xi}(\xi) 
			\end{equation*}
			\item Two-sided test:
			\begin{equation*}
				\tilde{\alpha} = 2 \cdot \min \left\{ P(\Xi \leq \xi), P(\Xi \geq \xi)\right\}
			\end{equation*}
		\end{itemize}
	\end{alertblock}
	%
	\newpage
	Examples:
	%
	\begin{table}%
		\begin{tabular}{l||c|c|c|c||c|c} \hline
			Type & $F_{\Xi}(\xi)$ & $\alpha$ & $\xi_{\mathrm{c,u}}$ & $\xi_{\mathrm{c,o}}$ & $\xi$ & $\tilde{\alpha} = P$-value\\ \hline
			\hline
			two-sided & SND & $0{,}05$ & $-1{,}960$ & $1{,}96$ & $0{,}00$ & $1{,}000$ \\ \hline
			two-sided & SND & $0{,}05$ & $-1{,}960$ & $1{,}96$ & $0{,}05$ & $0{,}960$ \\ \hline
			two-sided & SND & $0{,}05$ & $-1{,}960$ & $1{,}96$ & $1{,}00$ & $0{,}317$ \\ \hline
			two-sided & SND & $0{,}05$ & $-1{,}960$ & $1{,}96$ & $1{,}96$ & $0{,}050$ \\ \hline
			two-sided & SND & $0{,}05$ & $-1{,}960$ & $1{,}96$ & $3{,}00$ & $0{,}003$ \\ \hline
			\hline
			lower-sided & SND & $0{,}05$ & $-1{,}645$ & --- & $3{,}00$ & $0{,}999$ \\ \hline
			lower-sided & SND & $0{,}05$ & $-1{,}645$ & --- & $0{,}00$ & $0{,}500$ \\ \hline
			lower-sided & SND & $0{,}05$ & $-1{,}645$ & --- & $-1{,}00$ & $0{,}159$ \\ \hline
			lower-sided & SND & $0{,}05$ & $-1{,}645$ & --- & $-1{,}64$ & $0{,}050$ \\ \hline
			lower-sided & SND & $0{,}05$ & $-1{,}645$ & --- & $-3{,}00$ & $0{,}001$ \\ \hline
			\hline
			upper-sided & $\chi^2_{29}$ & $0{,}10$ & --- & $39{,}1$ & $5{,}00$ & $0{,}999$ \\ \hline
			upper-sided & $\chi^2_{29}$ & $0{,}10$ & --- & $39{,}1$ & $30{,}00$ & $0{,}414$ \\ \hline
			upper-sided & $\chi^2_{29}$ & $0{,}10$ & --- & $39{,}1$ & $39{,}00$ & $0{,}102$ \\ \hline
			upper-sided & $\chi^2_{29}$ & $0{,}10$ & --- & $39{,}1$ & $39{,}10$ & $0{,}100$ \\ \hline
			upper-sided & $\chi^2_{29}$ & $0{,}10$ & --- & $39{,}1$ & $50{,}00$ & $0{,}009$ \\ \hline
		\end{tabular}
	\end{table}
	\par\smallskip
	Usually, $\alpha$ is specified before the test, and the null hypothesis is then rejected if $\tilde{\alpha}$ is less than or equal to $\alpha$. The smaller $\tilde{\alpha}$ is compared to $\alpha$, the more \alert{significant} the rejection of ${\mathrm{ H}}_0$ is!
	%
\end{frame}

\end{document}



